\section{基于配对等待下界的双目标邻域搜索}\label{sec:third-innovation}

\subsection{从固定工期精修到交付与等待的权衡}

同一批次可以有两种交付要求。若交期优先，可固定阶段一所得的完工上界，再在该上界内压缩等待。若设备使用方愿意延迟部分产品回片以降低途中等待，则需要知道放宽工期能换来多少改善。沿用第三章配置D的可行域$\mathcal X_D$及事件计量口径，本文求解
\begin{equation}
 \min_{x\in\mathcal X_D}\bigl(A(x),B(x)\bigr),\qquad
 A=C_{\max},\quad B=L_{\rm sum}+10L_{\max}+25R .
 \label{eq:pareto-ab}
\end{equation}
$L_{\rm sum}$累计产品与辅助片扣除必要动作后的路线等待，$L_{\max}$为最大单项等待，$R$为同一留片链内相邻生产周期间隔之差的绝对值之和。两目标分别用秒与加权秒表示。若一份排程的两项目标都不大于另一份，且至少一项更小，则前者支配后者。非支配排程给出采购交期或工艺指标上限下可选择的资源安排。\par

给定允许的$B$上限$\varepsilon$，求解$\min A(x)$且$B(x)\le\varepsilon$，在工期相同的解中继续降低$B$。同时搜索两个单目标端点，并维护可行排程的非支配集合。这种上限法可以访问线性加权和遗漏的非支持点\citep{mavrotas2009}。本文将配对所强制的路线等待用于判断哪些离散改动值得进入子问题。

\subsection{以配对方式筛选邻域}

设$p$为产品与产品组成的PW对数，$d$为产品与真空辅助片组成的PW对数，$2p+d=N$。在标称时间参数下，双产品对共用一个对准器，第二片在第一片取离后还须经历放片、对准和取片。两片对准取片结束时刻至少相差$4+20+4=28$~s，却须同时装腔。对准后至共同装腔的三项额外等待中，较早一片累计至少28~s，至少一项达到$28/3$~s。产品与真空辅助片从上装载锁和真空存储区分别取片，两个站点之间一次取放与移动至少相隔$1+2=3$~s。共同装腔前与共同卸腔后各产生至少3~s等待。因此可用
\begin{equation}
 B(x)\ \ge\ 28p+6d+
 10\max\!\left\{\frac{28}{3}\mathbf{1}_{p>0},\,3\mathbf{1}_{d>0}\right\}
 \equiv \underline B(p,d).
 \label{eq:pareto-pair-bound}
\end{equation}
各对的等待属于不同晶圆路线，可以求和。最大项取所有对中成立的最大下界，节拍项非负。若$\underline B(p,d)>\varepsilon$，本次上限下即可跳过相应配对数量。换用其他设备时间时，应以该配置的对准、取放与移动时长重新计算下界。

例如，两片产品同对时，式~\eqref{eq:pareto-pair-bound}给出$B\ge28+280/3=364/3$。若设备要求$B\le100$，则无需再为这种配对调整腔室次序。两片产品各配一片辅助片时，下界降至$12+30=42$。这种结构通过筛选后，仍须求解并复验具体动作时刻，才能确定可交付工期。

\subsection{结构调整、重定时与验证}

从共同的可行排程出发，搜索改变产品配对、分腔、对内取放顺序和生产块竞争共享资源的次序。若两个父排程在某片产品的安排上不同，优先释放该片及其配对成员，并纳入相关腔室和相邻产品的资源次序。固定其余整数决策后，以线性规划调整动作时刻，并在小规模算例中用受限混合整数子问题修复联合邻域。每次替换均比较从逐片事件重算的两个目标。新增点删除其支配的旧点，档案保留可恢复的晶圆身份与动作表。学习选割沿用已有冻结策略，在带上限的完整模型分支中设置开关，单独评价其增量。

75片时，同一可行域生成的原模型含数百万个变量。大批次实现从共同可行排程提取两个Pjob内的有序生产块，保留Pjob取片顺序，在两个腔室之间重分配生产块、交换不同配方相邻块、拆分产品对或改变对内取放。领域搜索从已有非支配解中选择父排程，提高在设备内停留较久的产品被调整的概率。固定离散次序后，只为实际发生的动作建立重定时模型。复原排程时将动作起点取到$1/16$~s网格，重新检查事件、设备约束和指标上限。75片路径使用CPU上的重定时求解，完整原模型用于最终可行性复核。第~\ref{sec:pareto-evidence}~节分别报告通用小规模求解、两片专用完整性验证及大批次扩展的结果。
